%-------------------------
% William Ortega — Main resume (Mech / Fluids / Propulsion)
% Based on Jake Gutierrez's resume template (MIT license)
% Edits: skills on top, bold dates, larger name, 11pt, one-line headings
%-------------------------
\documentclass[letterpaper,11pt]{article}

\usepackage{latexsym}
\usepackage[empty]{fullpage}
\usepackage{titlesec}
\usepackage{marvosym}
\usepackage[usenames,dvipsnames]{color}
\usepackage{verbatim}
\usepackage{enumitem}
\usepackage[hidelinks]{hyperref}
\usepackage{fancyhdr}
\usepackage[english]{babel}
\usepackage{tabularx}
\input{glyphtounicode}

\pagestyle{fancy}
\fancyhf{}
\fancyfoot{}
\renewcommand{\headrulewidth}{0pt}
\renewcommand{\footrulewidth}{0pt}

% Margins (if it spills onto page 2, shrink these a bit first)
\addtolength{\oddsidemargin}{-0.6in}
\addtolength{\evensidemargin}{-0.6in}
\addtolength{\textwidth}{1.2in}
\addtolength{\topmargin}{-0.7in}
\addtolength{\textheight}{1.4in}

\urlstyle{same}
\raggedbottom
\raggedright
\setlength{\tabcolsep}{0in}

\titleformat{\section}{
  \vspace{-6pt}\scshape\raggedright\large\bfseries
}{}{0em}{}[\color{black}\titlerule \vspace{-6pt}]

\pdfgentounicode=1

%----------- Custom commands -----------
\newcommand{\resumeItem}[1]{\item\small{#1 \vspace{-2.5pt}}}

% Two-line heading (used for Education)
\newcommand{\resumeSubheading}[4]{
  \vspace{-2pt}\item
    \begin{tabular*}{\textwidth}[t]{l@{\extracolsep{\fill}}r}
      \textbf{#1} & \textbf{\small #2} \\
      \textit{\small#3} & \textit{\small #4} \\
    \end{tabular*}\vspace{-7pt}
}

% One-line heading: Title ............ Dates (bold)
\newcommand{\resumeRole}[2]{
  \vspace{-2pt}\item
    \begin{tabular*}{\textwidth}[t]{l@{\extracolsep{\fill}}r}
      \textbf{#1} & \textbf{\small #2} \\
    \end{tabular*}\vspace{-7pt}
}

\renewcommand\labelitemii{$\vcenter{\hbox{\tiny$\bullet$}}$}
\newcommand{\resumeSubHeadingListStart}{\begin{itemize}[leftmargin=0.15in, label={}]}
\newcommand{\resumeSubHeadingListEnd}{\end{itemize}}
\newcommand{\resumeItemListStart}{\begin{itemize}}
\newcommand{\resumeItemListEnd}{\end{itemize}\vspace{-7pt}}

%=====================================================
\begin{document}

%----------HEADING----------
\begin{center}
    {\fontsize{26}{30}\selectfont \textbf{William Ortega}} \\ \vspace{3pt}
    \small Ontario, CA $|$ 909-228-3985 $|$ wortega@g.hmc.edu $|$
    \href{https://william-ort.github.io/}{\underline{william-ort.github.io}}
\end{center}

%-----------EDUCATION-----------
\section{Education}
  \resumeSubHeadingListStart
    \resumeSubheading
      {Harvey Mudd College}{Expected May 2028}
      {Bachelor of Science in Engineering, GPA: 3.4}{Claremont, CA}
      \resumeItemListStart
        \resumeItem{\textbf{Relevant Coursework:} Fluid Mechanics, Thermodynamics, Statics, Experimental Engineering, Design \& Manufacturing, Signals \& Systems, Differential Equations \& Fourier Analysis, AI-Amplified Prototyping}
      \resumeItemListEnd
  \resumeSubHeadingListEnd

%-----------TECHNICAL SKILLS-----------
\section{Technical Skills}
 \begin{itemize}[leftmargin=0.15in, label={}]
    \small{\item{
     \textbf{CAD \& Simulation}{: SolidWorks, Ansys (CFD), FEA, COMSOL, OpenRocket} \\
     \textbf{Programming}{: MATLAB, Python, C++, LabVIEW, SystemVerilog} \\
     \textbf{Fabrication}{: 2-Axis Machining, 3D Printing, Laser \& Waterjet Cutting, GD\&T, DFM/DFA} \\
     \textbf{Languages}{: English (Native), Spanish (Native)}
    }}
 \end{itemize}

%-----------EXPERIENCE-----------
\section{Engineering Experience}
  \resumeSubHeadingListStart

    \resumeRole
      {Undergraduate Research Assistant, Flow Imaging Lab, Harvey Mudd College}{Jan. 2026 -- Present}
      \resumeItemListStart
        \resumeItem{Demonstrated higher thrust coefficients for submerged horizontal fins than surface-crossing fins (peak mean $C_T \approx 4.9$) by designing a horizontal-fin test fixture and running 100+ trials across fin thickness, amplitude, and frequency}
        \resumeItem{Characterized how aspect ratio, flapping frequency, and submersion depth affect thrust across 100+ trials with an ATI Mini40 six-axis load cell and NI DAQ, holding frequency within $\pm$0.02 Hz}
        \resumeItem{Redesigned the lab's MATLAB thrust-processing pipeline for horizontal fins by remapping the thrust axis and rewriting the in-air inertial correction, enabling direct thrust-coefficient comparison with vertical fins}
      \resumeItemListEnd

    \resumeRole
      {Propulsion Subteam Member, Mudd Amateur Rocketry Club (MARC)}{Sept. 2026 -- Present}
      \resumeItemListStart
        \resumeItem{Sized canards for the team's competition rocket with OpenRocket's optimizer, identifying dimensions that kept the stability margin above 1.2 calibers while preserving control authority}
        \resumeItem{Set up Ansys CFD simulations of the canards and fin tabs to build an aerodynamic lookup table vs. angle of attack for the rocket's flight controller}
      \resumeItemListEnd

  \resumeSubHeadingListEnd

%-----------PROJECTS-----------
\section{Engineering Projects}
    \resumeSubHeadingListStart

      \resumeRole
          {Autonomous Underwater Vehicle (Team of 4)}{Jan. 2026 -- May 2026}
          \resumeItemListStart
            \resumeItem{Built and deployed a PVC-frame AUV with proportional depth control that profiled light, pressure, and temperature off Dana Point, CA, measuring $>$80\% surface-light loss by 3 m via Beer--Lambert analysis in MATLAB}
            \resumeItem{Designed and soldered a sensor/DAQ board integrating a TCS34725 RGB sensor (I\textsuperscript{2}C), MPX5700 pressure sensor, and NTC thermistor on a Teensy 4.1, logging depth-resolved data to microSD during motors-off holds to cut noise}
            \resumeItem{Designed a two-stage MCP6002 op-amp circuit (gain $\approx$ 6.87) that scaled the pressure signal into the 0.3--3.3 V ADC window to prevent clipping, and calibrated voltage-to-depth in saltwater with a linear fit}
          \resumeItemListEnd

      \resumeRole
          {Mobile Lectern, Client Design Project (Team of 4)}{Sept. 2025 -- Nov. 2025}
          \resumeItemListStart
            \resumeItem{Designed and built a 3 lb, \$24 collapsible lectern for a professor client that sets up in 1.6 s and holds over 12 lb}
            \resumeItem{Created the SolidWorks CAD model, drawings, and BOM, and laser-cut eco-friendly material for the final prototype}
          \resumeItemListEnd

      \resumeRole
          {Machinist's Hammer}{Sept. 2025 -- Nov. 2025}
          \resumeItemListStart
            \resumeItem{Calibrated and operated manual mills and lathes to machine a hammer from raw steel and nylon to $\pm$0.001 in.}
            \resumeItem{Applied drilling, tapping, heat treatment, and tempering while following technical drawings and a process router}
          \resumeItemListEnd

    \resumeSubHeadingListEnd

%-----------LEADERSHIP-----------
\section{Leadership \& Activities}
  \resumeSubHeadingListStart

    \resumeRole
      {Co-Founder \& Co-President, SHPE Harvey Mudd Chapter}{Jan. 2026 -- Present}
      \resumeItemListStart
        \resumeItem{Grew chapter engagement to 52 students through recruiting, outreach, and professional-development programming}
        \resumeItem{Secured \$10K in institutional funding to send 16 members to the 2026 SHPE National Convention}
        \resumeItem{Planned 3 professional-development workshops and launched company outreach to build a sponsorship pipeline}
      \resumeItemListEnd

    \resumeRole
      {President, Societal Advancement of Latinx (SALSA), Harvey Mudd College}{Jan. 2025 -- May 2026}
      \resumeItemListStart
        \resumeItem{Led and planned 5+ cultural and social events each semester to build community and belonging}
        \resumeItem{Expanded collaboration with five Latinx organizations across the Claremont Colleges}
      \resumeItemListEnd

  \resumeSubHeadingListEnd

\end{document}
